\documentclass{article}
\usepackage[utf8]{inputenc}
\usepackage[spanish]{babel}
\usepackage{listings}
\usepackage{xcolor}
\usepackage{geometry}
\geometry{a4paper, margin=1in}

\% Configuración para código C++
\definecolor{codegreen}{rgb}{0,0.6,0}        
\definecolor{codegray}{rgb}{0.5,0.5,0.5}     
\definecolor{codepurple}{rgb}{0.58,0,0.82}   
\definecolor{backcolour}{rgb}{0.95,0.95,0.92}

\lstset{
    backgroundcolor=\color{backcolour},      
    commentstyle=\color{codegreen},
    keywordstyle=\color{magenta},
    numberstyle=\tiny\color{codegray},       
    stringstyle=\color{codepurple},
    basicstyle=\ttfamily\small,
    breakatwhitespace=false,
    breaklines=true,
    captionpos=b,
    keepspaces=true,
    numbers=left,
    numbersep=5pt,
    showspaces=false,
    showstringspaces=false,
    showtabs=false,
    tabsize=2,
    language=C++,
    inputencoding=utf8,
    extendedchars=true,
    literate={á}{{\'a}}1 {é}{{\'e}}1 {í}{{\'i}}1 {ó}{{\'o}}1 {ú}{{\'u}}1 {ñ}{{\~n}}1 {¿}{{\textquestiondown}}1 {¡}{{\textexclamdown}}1
}

\% Macros para las secciones de los apuntes
\newcommand{\quees}[1]{
    \section*{🤔 ¿QUÉ ES?}
    \hrulefill \\
    \#1
}

\newcommand{\dichodeotraforma}[1]{
    \section*{💬 DICHO DE OTRA FORMA}
    \hrulefill \\
    \#1
}

\newcommand{\diagrama}[1]{
    \section*{🎨 DIAGRAMA}
    \hrulefill \\
    \#1
}

\newcommand{\comofunciona}[1]{
    \section*{⚙️ ¿CÓMO FUNCIONA?}
    \hrulefill \\
    \#1
}

\newcommand{\complejidad}[1]{
    \section*{📱 COMPLEJIDAD (Big-O)}
    \hrulefill \\
    \#1
}

\newcommand{\entrevistas}[1]{
    \section*{🎯 EN ENTREVISTAS}
    \hrulefill \\
    \#1
}

\newcommand{\resumen}[1]{
    \section*{📋 RESUMEN RÁPIDO}
    \hrulefill \\
    \#1
}

\begin{document}

\title{Linked Lists (Listas Enlazadas)}
\author{Aldo Zepeda Montoya}
\date{\today}
\maketitle

\subsection*{CATEGORÍA: 01\_Estructuras\_Lineales}

\quees{
Una cacería del tesoro 🗺️. Cada pista (node) te dice dónde está la 
SIGUIENTE pista. 

No sabes dónde está la pista \#5 sin seguir la cadena desde el inicio: 
pista 1 → pista 2 → pista 3 → pista 4 → pista 5. 

A diferencia de los arrays (vagones pegaditos), los nodos de una 
Linked List pueden estar en cualquier parte de la memoria. Lo único
que los mantiene unidos es que cada uno sabe la dirección (pointer)
del siguiente.
}

\dichodeotraforma{
Una Linked List es una estructura lineal donde los elementos no se
guardan en posiciones contiguas de memoria. Cada elemento es un objeto
llamado "Node", que contiene:
1. El valor (data).
2. Un puntero (next) al siguiente nodo.

Esto permite insertar y eliminar elementos en los extremos en O(1) 
sin tener que desplazar a nadie, pero sacrificas el acceso aleatorio
(no puedes ir directo al índice `i`).
}

\diagrama{\begin{verbatim}
Singly Linked List:
HEAD
 ↓
[3|•]──→[7|•]──→[1|•]──→[9|NULL]
 node0   node1   node2   node3

Doubly Linked List:
HEAD                      TAIL
 ↓                         ↓
[NULL|3|•]⇄[•|7|•]⇄[•|1|•]⇄[•|9|NULL]

Insertar 5 después del node1 (valor 7):
ANTES:   [3]──→[7]──→[1]──→[9]
                ↓
DESPUÉS: [3]──→[7]──→[5]──→[1]──→[9]   ← O(1) si ya tienes el puntero
\end{verbatim}}

\begin{lstlisting}[language=C++]
1. Acceso/Búsqueda: Tienes que empezar en el HEAD e ir saltando de 
   nodo en nodo. O(n).
2. Inserción:
   - Al Head: Creas nodo, apuntas su `next` al viejo head, actualizas
     head. O(1).
   - En medio: Llegas a la posición, ajustas los punteros para incluir
     el nuevo nodo.
3. Eliminación:
   - Del Head: Mueves el head al `head->next`. O(1).
   - En medio: Ajustas el puntero del nodo anterior para que "salte"
     al nodo que quieres borrar.

📝 PSEUDOCÓDIGO
──────────────────────────────────────────────────────────────

// Insertar al inicio
función insertarAlInicio(valor):
    nuevoNodo = crear Nodo(valor)
    nuevoNodo.next = head
    head = nuevoNodo

// Buscar valor
función buscar(valor):
    actual = head
    MIENTRAS actual != NULL:
        SI actual.valor == valor: return VERDADERO
        actual = actual.next
    return FALSO

💻 CÓDIGO C++
──────────────────────────────────────────────────────────────

#include <iostream>
using namespace std;

// Definición del Nodo
struct Node {
    int data;
    Node* next;
    Node(int val) : data(val), next(nullptr) {} // Constructor
};

class LinkedList {
public:
    Node* head;
    LinkedList() : head(nullptr) {}

    // Insertar al inicio -> O(1)
    void insertHead(int val) {
        Node* newNode = new Node(val);
        newNode->next = head;
        head = newNode;
    }

    // Insertar al final -> O(n) (sin tail pointer)
    void insertTail(int val) {
        Node* newNode = new Node(val);
        if (!head) {
            head = newNode;
            return;
        }
        Node* temp = head;
        while (temp->next) temp = temp->next; // Recorrer hasta el final
        temp->next = newNode;
    }

    // Buscar -> O(n)
    bool search(int val) {
        Node* temp = head;
        while (temp) {
            if (temp->data == val) return true;
            temp = temp->next;
        }
        return false;
    }

    // Imprimir lista
    void print() {
        Node* temp = head;
        while (temp) {
            cout << temp->data << " -> ";
            temp = temp->next;
        }
        cout << "NULL" << endl;
    }
};

int main() {
    LinkedList list;
    list.insertHead(10);
    list.insertHead(20);
    list.insertTail(5);
    
    list.print(); // 20 -> 10 -> 5 -> NULL
    
    cout << "¿Está el 10? " << (list.search(10) ? "Sí" : "No") << endl;
    
    return 0;
}

⏱️ COMPLEJIDAD (Big-O)
──────────────────────────────────────────────────────────────

  Operación        │ Tiempo  │ vs Array
  ─────────────────┼─────────┼──────────────────────────────
  Access por índice│ O(n)    │ Array: O(1) ← array gana
  Search           │ O(n)    │ Igual
  Insert head      │ O(1)    │ Array: O(n) ← linked list gana
  Insert tail      │ O(1)*   │ Array: O(1) amort. → Igual
  Insert medio     │ O(n)    │ Igual (pero no desplaza datos)
  Delete head      │ O(1)    │ Array: O(n) ← linked list gana
  *con tail pointer
\end{lstlisting}

\entrevistas{
✅ Patrón "Fast \& Slow Pointers" (Tortoise and Hare):
   - Sirve para detectar ciclos (si se encuentran, hay ciclo).
   - Sirve para encontrar el medio (cuando fast llega al final, slow 
     está en el medio).

💡 Dummy Node: Crear un nodo falso al inicio ayuda a simplificar casos
   borde (como insertar en una lista vacía o borrar el head).

⚠️ Memory Management: En C++, recuerda que cada `new` necesita un `delete`
   para evitar memory leaks (fugas de memoria).

🔑 Pregunta clave: "¿Es singly o doubly?". Si es doubly, puedes ir hacia
   atrás, lo que facilita borrar un nodo si ya tienes el puntero.
}

\resumen{
• Cadena de nodos dispersos en memoria.
• Acceso secuencial (no puedes saltar).
• Excelente para insertar/eliminar al inicio O(1).
• No desperdicia memoria (solo usa lo que necesita).
• Invertir una Linked List es un problema clásico (usar 3 punteros).
• Usa más memoria por elemento que un array (por el puntero next).
════════════════════════════════════════════════════════════════
}

\end{document}